% -*- mode:LaTex; mode:visual-line; mode:flyspell; fill-column:75-*-

% Special indentation for abstract.
\setlength{\parskip}{1em}
\setlength{\parindent}{0em}

\noindent
United States apple growers lose more than \$100 million a year to fire blight,
a bacterial disease caused by \textit{Erwinia amylovora} that can kill a young,
high-density planting in a single season.
The disease is best fought during the dormant season: the infected tissue that
survives winter on the tree is what seeds the following spring's outbreak, so
removing that tissue before spring growth begins is the most effective way to
stop the disease from spreading.
Finding it still depends almost entirely on manual scouting, a slow and
subjective process that does not scale to the multi-acre plantings of modern
commercial operations and is increasingly constrained by a shrinking
agricultural workforce.

This thesis develops a robotic system capable of performing that inspection
autonomously.
Because the visual cues that distinguish infected tissue are subtle, easily
obscured by occlusion and variable natural lighting, and reliably resolved only
at close range, the system relies on active perception: a manipulator positions
a camera to acquire discriminative, task-relevant views of the canopy.
We built a multi-modal dataset of dormant apple trees, the first to pair dense
near-infrared imagery with flash-illuminated stereo RGB for this task, and
trained instance-segmentation detectors to recognize disease symptoms across
both modalities.
We then introduced a confidence-aware semantic mapping method that fuses these
per-view detections into a persistent 3D representation of disease confidence
across the canopy, and a next-best-view planner that actively selects
viewpoints to refine the map's least confident, most contested predictions.
Finally, we integrated the full pipeline onto Erwin, a mobile manipulator built
on an Amiga base and an xArm6 arm, and deployed it on real dormant apple trees
in a working orchard.

Validated in simulation and in the field on 12 trees at the Penn State Fruit
Research and Extension Center under a genuine train--test domain gap, the
semantic planner reaches detection quality comparable to a complete fixed scan
in roughly half the viewpoints and half the inspection time.
This result demonstrates the central promise of active perception for orchard
disease mapping: by deciding where to look next rather than only looking
harder, a robot can build maps of disease that are both more accurate and
faster to produce than one that follows a fixed scanning routine.
